\section{引言}
可复用技能将说明、脚本与资源组织为能力单元，技能网络再通过检索和编排连接任务。
但文档完整、模型评分高、节点更多或检索相似度高，均不足以推出指定任务能够完成。
本研究关注两个直接问题：怎样得到可核验的技能质量判断及能力边界；怎样利用这些证据让技能网络在明确预算内完成更多任务。
评判依据应是可回放反证、独立完整产物和资源开销，而不是网络规模或模型自报成功。

SkillNet将技能创建、评价与连接整合为网络框架\citep{P01}；
SkillCT提出静态、动态与人类辅助画像\citep{P43}；SkillLearnBench区分技能文本、轨迹与任务结果\citep{P44}。
SkillsBench与ASSAY强调执行对照和条件化效应\citep{P02,P03}，
路由及多样性研究进一步关注大候选池与互补选择\citep{P04,P05}。
因此本研究不以新增评分维度、规则数量或确定性工具本身作为算法创新。
实质问题是：哪些证据能支持哪种决定，何种组合在什么条件下产生可重复增量，何时强而简单的对照已足够。

在SkillNet-demo基础上，本文交付两个可运行成果。
第一项是证据约束技能判别器，将固定版本的资源、机械接口反证与可信产物回执连接，
输出需要修复、被测范围通过、证据不足和超出已测范围。
第二项是证据条件化路径方法ECP，在有AND前提和预算约束的技能网络中，按版本及输入条件匹配开发证据，
保留互补的完整路径，并在新任务上重新验收。
前者不把静态一致当作运行认证；后者不把开发成功当作未来成功保证。

论文回答三个可检验问题：
\begin{itemize}
\item RQ1：机械反证和行为回执能否区分具有相似高文档分的不同实现，已测能力边界如何表达？
\item RQ2：条件证据与兼容路径能否在相同候选、预算和执行器下，比原生最短路径、全局证据路径及人工固定流程完成更多任务？
\item RQ3：增量来自条件化、集合互补还是版本绑定，取证成本和未见压力如何限制结论？
\end{itemize}

前一阶段实证为600例判别诊断、288次自建程序输出，以及公开数据派生的18节点网络执行。
网络开发使用两份数据，主留出使用另两份数据和128个构造任务；十种策略、三档预算及压力结果完整报告。
相对原生元数据路径，ECP得到明确条件收益；相对人工与全局风险路径，成功增量主要限于紧预算。
合并新数值精度与版本变化后出现12/24的失败边界，前期取证也使本批总调用开销高于人工流程。
这些正负结果共同界定当前成果，不能推广为1,881个自然SkillNet节点的普遍优化或领域质量认证。

作为研究动机与测量参照，论文同时保留200个自然包审查、381个接口干预、54次模型评价、
300次五臂文本产物及1,152条既见网络回放。
这些历史实验不与新留出混合计算样本。本轮另增加社区原始CLI隔离执行、最小修复与外部真实轨迹评审，分别披露付费调用中断及样本限制。
可核验性来自公开实现、固定输入与反例，而非引用数量或模板的期刊外观。

本轮新增结果进一步检验而非只扩展上述构造成功：355次社区原始/修复工具执行证实两处具体失败可修复，
但外部86例完整配对的模型评审没有可靠优势。自然工具组合回放也显示修复后的简单路径足够，
不能将节点能力改进全归于复杂编排器。这些结果把方法的有效部分与尚未成立的外推分开。
